\section{Introduction}
Automatic fake-news detectors are routinely reported with accuracies above 99\% on public benchmarks, yet the same models often degrade sharply when evaluated on a different dataset~\cite{janicka2019cross,hoy2022generalisability,delloglio2026comparison}. A common explanation is that benchmarks contain artifacts: properties that correlate with the label in one collection but say little about veracity in another, for example publisher markers, document length, or the entities that happened to be in the news when the data were collected. This is an instance of shortcut learning~\cite{geirhos2020shortcut}, well documented for natural language inference~\cite{gururangan2018annotation,poliak2018hypothesis,mccoy2019right} and increasingly for misinformation detection~\cite{verma2026accuracy,wan2025truth,xiao2023sok,haroon2026benchmarks}. However, showing that a cue correlates with the label is not evidence that removing it would improve transfer: a gap that shrinks after text is removed may reflect the perturbation itself rather than the targeted cue. Controlled interventions, with size-matched random deletion as a matched perturbation control, are therefore needed.

Prior work documents transfer failures, entity bias and sensitivity to topical or stylistic cues~\cite{delloglio2026comparison,zhu2022endef,murayama2021diachronic,wan2025truth}. These studies motivate three questions: (1) how much of the measured cross-dataset gap is removed by the selected interventions, (2) whether masking gains exceed those of equally sized random deletion on the same articles, and (3) how responses vary across dataset pairs and model families. Our contribution is the joint diagnostic design addressing these questions, combining a TF-IDF factorial study with document-paired masking controls and a focused two-transformer validation. It does not depend on claiming that earlier work never used random perturbations or investigated artifacts.

We make four contributions.
\begin{enumerate}
\item \textbf{A controlled attribution framework.} A factorial design over source markers, length and entities, with size-matched random-deletion controls (including controls for the length-matched combinations), equal-size class-balanced samples, document-paired masking comparisons, group-aware splits, overlap-aware datasets and per-pair reporting.
\item \textbf{A dataset audit and construction.} Under our normalized matching every cleaned ISOT and McIntire document has a WELFake match, ISOT carries about 13\% exact duplicates, and the WELFake label direction is described inconsistently; we resolve it empirically and build an overlap-trimmed residual.
\item \textbf{Findings for a linear model.} In the paired validation, 26.2\% of the TF-IDF cross-dataset AUC gap is removed by the three interventions, almost entirely through length and entity cues; archived exploratory results suggest transferable web-boilerplate signal; entity masking beats its size-matched random control, also after length matching; and the effects differ strongly by dataset pair.
\item \textbf{A model-dependence result.} In document-paired validation, entity masking performs worse than matched random deletion for DistilBERT, while neither transformer has an established positive pooled S+L+E effect beyond its length-matched random control. Responses vary across models and dataset pairs; a large raw gap reduction does not establish a cue-specific benefit.
\end{enumerate}
All claims are limited to three English news corpora from 2016--2017 and three models (a linear lexical model and two fine-tuned transformers).

\subsection{Objectives}
O1 establishes the audited analysis corpora. O2 addresses the first research question, O3 the second, and O4 the third.
\begin{enumerate}
\item[O1] Audit three public corpora for overlap, duplication and label direction, and construct a defensible set of datasets.
\item[O2] Quantify, for a TF-IDF logistic-regression detector, how much of the cross-dataset gap is removed by source-marker masking, length matching and entity masking, alone and combined.
\item[O3] Compare source-marker and entity masking, alone and combined with length matching, against deletion of an equal number of random words; compare length matching alone against unconstrained sampling.
\item[O4] Report how the effects vary across ordered dataset pairs, and whether they carry over to fine-tuned transformers (DistilBERT and RoBERTa).
\end{enumerate}



